%
%%%%

%%%   Самарский государственный технический университет

%%%   Кафедра прикладной математики и информатики

%%%

%%%   Преамбула для статьи в журнал "Вестник Самарского государственного

%%%   технического университета. Серия: «Физико-математические науки»"

%%%   Ориентирована под MikTEx 2.4 for Win32

%%%

%%%   Хотя сам журнал собирается под GNU/Linux на дистрибутиве TeXLive



\documentclass[11pt,a4paper,twoside]{article}



\usepackage[cp1251]{inputenc}

\usepackage[T2A]{fontenc}

\usepackage[english,russian]{babel}

\usepackage{samgtubib}

\usepackage[dvips]{graphicx}

\usepackage[multiple]{footmisc}



\usepackage{samgtu}

\renewcommand{\VestnikYear}{2009} % Год издания Вестника

\renewcommand{\VestnikNumber}{2\,(19)} % Номер Вестника

\renewcommand{\VestnikMonth}{Ноябрь} % Месяц выхода Вестника

\renewcommand{\VestnikData}{01.11.09} % Дата подписания Вестника в печать

\renewcommand{\VestnikUPL}{26,64} % Усл. печ. л.

\renewcommand{\VestnikUIL}{25,73} % Уч.-изд. л.

\renewcommand{\VestnikRegNumber}{479/09} % Регистрационный номер

\renewcommand{\VestnikOrder}{994} % Номер заказа






%
%\begin{document}

\mathtoolsset{showmanualtags} %

\renewcommand{\firstpage}{305}
\renewcommand{\lastpage}{314}

\udk{004.38:530.145}

\msc{81P99; 81V80, 94A50}



\email{E-mail}{watanabe@is.noda.tus.ac.jp}



\engtitle{Note on complexity of quantum transmission processes}

\engauthor{N.~Watanabe}{N.~W\,a\,t\,a\,n\,a\,b\,e}



\enginstitute{Science University of Tokyo,    Noda City,  Chiba 278-8510, Japan.
}

\otherenginf{
\fullauthor{Noboru Watanabe} \position{Professor} \dept{Dept. of Information Sciences} 
}




\engabstract{In 1989, Ohya propose a new concept, so-called Information Dynamics (ID), 
to investigate complex systems according to two kinds of view points. One is the dynamics 
of state change and another is measure of complexity.  In ID, two complexities  $  C^{S} $   and  $  T^{S} $   
are introduced.   $  C^{S} $   is a measure for complexity of system itself, and  $  T^{S} $   is a measure 
for dynamical change of states, which is called a transmitted complexity.  An example of these 
complexities of ID is entropy for information transmission processes.  The study of complexity 
is strongly related to the study of entropy theory for classical and quantum systems. 
The quantum entropy was introduced by von Neumann around 1932, which describes 
the amount of information of the quantum state itself. It was extended by Ohya for 
C*-systems before CNT entropy. The quantum relative entropy was first defined by Umegaki for 
 $  \sigma $-finite von Neumann algebras, which was extended by Araki and Uhlmann for general 
von Neumann algebras and *-algebras, respectively. By introducing a new notion, the 
so-called compound state, in 1983 Ohya succeeded to formulate the mutual entropy 
in a complete quantum mechanical system (i.e., input state, output state and channel are 
all quantum mechanical) describing the amount of information correctly transmitted through 
the quantum channel.
In this paper, we briefly review the entropic complexities for classical and 
quantum systems. We introduce some complexities 
by means of entropy functionals in order to treat the transmission processes consistently. 
We apply the general frames of quantum communication to the Gaussian communication processes.  
Finally, we discuss about a construction of compound states including quantum correlations.}

\engkeywords{quantum communication channel, von Neumann entropy, $\mathcal{S}$-mixing entropy, Ohya mutual entropy, C*-system}

\date{14/XI/2012}
\revisiondate{21/I/2013}




\makeengtitlef

\smallskip


\textbf{1. Introduction}

In \cite{nv:-O3}, Ohya introduced Information Dynamics (ID) synthesizing
dynamics of state change and complexity of state. Based on ID, one can study
various problems of physics and other fields. Channel and two complexities
are key concepts of ID.

Let us briefly review ID  for quantum communication processes.

Let  $  \mathcal{H}_{k} $    $  (k=1, 2) $   be complex separable Hilbert spaces.   We
denote the set of all bounded linear operators on  $  \mathcal{H}_{k} $   by  $  %
\mathbf{B}(\mathcal{H}_{k}) $    $  (k=1, 2) $   and we express the set of all density
operators on  $  \mathcal{H}_{k} $   by  $  \mathfrak{S(}\mathcal{H}_{k}) $    $  (k=1, 2) $.
  Let  $  \left( \mathbf{B}(\mathcal{H}_{k}),\mathfrak{S(}\mathcal{H}%
_{k})\right)  $    $  (k=1, 2) $   be input  $  (k=1) $   and output  $  (k=2) $   quantum
systems, respectively.

\smallskip

\textbf{1.1. Quantum Channels}

A mapping from  $  \mathfrak{S(}\mathcal{H}_{1}) $   to  $  \mathfrak{S(}\mathcal{H}%
_{2}) $   is called a quantum channel  $  \Lambda ^{\ast } $.

\begin{enumerate}
\item[(1)]  $  \Lambda ^{\ast } $  \textit{  }is called a linear channel if  $  %
\Lambda ^{\ast } $   satisfies the affine property such as  $  \Lambda ^{\ast
}(\sum_{k}\lambda _{k}\rho _{k})=\sum_{k}\lambda _{k}\Lambda ^{\ast }\left(
\rho _{k}\right)  $   for any  $  \rho _{k}\in \mathfrak{S(}\mathcal{H}_{1}) $   and
any nonnegative number  $  \lambda _{k}\in \lbrack 0,1] $   with  $  \sum_{k}\lambda
_{k}=1 $.
\end{enumerate}

For the quantum channel  $  \Lambda ^{\ast } $,   the dual map  $  \Lambda  $   of  $  %
\Lambda ^{\ast } $   is defined by 
\[
\mathop{\rm tr}\Lambda ^{\ast }(\rho )B=\mathop{\rm tr}\rho \Lambda (B),\quad \forall \rho \in 
\mathfrak{S(}\mathcal{H}_{1}),    \forall B\in \mathbf{B}(\mathcal{H}_{2}). 
\]

\begin{enumerate}
\item[(2)]  $  \Lambda ^{\ast } $    is called a completely positive (CP) channel
if  $  \Lambda ^{\ast } $    is linear channel and its dual map  $  \Lambda :\mathbf{B%
}(\mathcal{H}_{2})\rightarrow \mathbf{B}(\mathcal{H}_{1}) $    of  $  \Lambda
^{\ast } $   holds 
\[
\left\langle x,\sum_{i,j=1}^{n}A_{i}^{\ast }\Lambda (\overline{A}_{i}^{\ast }%
\overline{A}_{j})A_{j}\,x\right\rangle \geq 0\quad \left( \forall x\in 
\mathcal{H}_{1}\right) 
\]%
for any  $  n\in \mathbf{N} $,   any  $  \{\overline{A}_{i}\}\subset \mathbf{B}(%
\mathcal{H}_{2}) $   and any  $  \{A_{i}\}\subset \mathbf{B}(\mathcal{H}_{1}) $.
\end{enumerate}

One can describe almost all physical transform of states by using the CP
channel [\citen{nv:-O1}--\citen{nv:-Ing}].

\smallskip

\textbf{1.2. Quantum Communication Channel}

Here we explain the quantum communication channels as an example of the
quantum channels.   

In order to consider influence of the environment such as noise and loss, we
suppose  $  \mathcal{K}_{1} $   and  $  \mathcal{K}_{2} $   to be complex separable
Hilbert spaces of noise and loss systems, respectively. Quantum channel of
quantum communication process with noise and loss was discussed by \cite{nv:-O1,nv:-OW3}.

\smallskip

\textbf{1.3. Noisy quantum channel and Generalized Beam Splitter}

\chead[\fancyplain{}{\scriptsize \sl N.~W\,a\,t\,a\,n\,a\,b\,e}]{\fancyplain{}{\scriptsize \sl
Note on complexity of quantum transmission processes}}


For an input state  $  \rho  $   in  $  \mathfrak{S(}\mathcal{H}_{1}\mathcal{)} $   and
a noise state  $  \xi \in \mathfrak{S}\left( \mathcal{K}_{1}\right)  $,   Ohya and
NW defined in \cite{nv:-OW3} a generalized beam splitting  $  \Pi ^{\ast } $   by 
\[
\Pi ^{\ast }(\rho \otimes \xi )\equiv V\left( \rho \otimes \xi \right)
V^{\ast }, 
\]%
where  $  V $   is a linear mapping from  $  \mathcal{H}_{1}\otimes \mathcal{K}_{1} $  
to  $  \mathcal{H}_{2}\otimes \mathcal{K}_{2} $   given by for the  $  n_{1} $,    $  m_{1} $,    $  j $,    $  \left( n_{1}+m_{1}-j\right)  $   photon number state vectors  $  \left\vert n_{1}\right\rangle \in \mathcal{H}_{1} $,    $  |m_{1}\rangle \in \mathcal{
K}_{1} $,    $  |j\rangle \in \mathcal{H}_{2} $,    $  |n_{1}+m_{1}-j\rangle \in \mathcal{K}_{2} $   
\[
V(|n_{1}\rangle \otimes |m_{1}\rangle
)=\sum_{j}^{n_{1}+m_{1}}C_{j}^{n_{1},m_{1}}|j\rangle \otimes
|n_{1}+m_{1}-j\rangle 
\]%
and 
\begin{equation}
\tag{1}
\begin{array}{rcl}
\displaystyle
C_{j}^{n_{1},m_{1}} 
&=&\displaystyle\sum_{r=L}^{K}(-1)^{n_{1}+j-r}{\frac{\sqrt{n_{1}!m_{1}!j!(n_{1}+m_{1}-j)!}%
}{{r!(n_{1}-j)!(j-r)!(m_{1}-j+r)!}}}  \\
&&\times \alpha ^{m_{1}-j+2r}\left(-\bar{\beta}\right) ^{n_{1}+j-2r} 
\end{array}
\end{equation}
 $  \alpha  $   and  $  \beta  $   are complex numbers satisfying  $  \left\vert \alpha
\right\vert ^{2}+\left\vert \beta \right\vert ^{2}=1 $.  $  K $   and  $  L $   are
constants given by  $  K=\min \{n_{1},j\} $,    $  L=\max \{m_{1}-j, 0\} $. For the
coherent input state  $  \rho =\left\vert \theta \right\rangle \left\langle
\theta \right\vert \otimes \left\vert \kappa \right\rangle \left\langle
\kappa \right\vert \in \mathfrak{S}\left( \mathcal{H}_{1}{\otimes \mathcal{K}%
}_{1}\right),  $   the output state of  $  \Pi ^{\ast } $   is obtained by 
\begin{eqnarray*}
\Pi ^{\ast }\left( \left\vert \theta \right\rangle \left\langle \theta
\right\vert \otimes \left\vert \kappa \right\rangle \left\langle \kappa
\right\vert \right) &=&\left\vert \alpha \theta +\beta \kappa \right\rangle
\left\langle \alpha \theta +\beta \kappa \right\vert \\
&&\otimes \left\vert-\bar{\beta}\theta +\bar{\alpha}\kappa \right\rangle
\left\langle-\bar{\beta}\theta +\bar{\alpha}\kappa \right\vert.
\end{eqnarray*}%
By using  $  \Pi ^{\ast } $,   Ohya and NW introduced in \cite{nv:-OW3} the noisy
quantum channel  $  \Lambda ^{\ast } $   with a fixed noise state  $  \xi \in 
\mathfrak{S}\left( \mathcal{K}_{1}\right)  $   defined by 
\begin{equation}
\Lambda ^{\ast }(\rho )\equiv \mathop{\rm tr}\nolimits_{\mathcal{K}_{2}}\Pi ^{\ast }(\rho \otimes
\xi )=\mathop{\rm tr}\nolimits_{\mathcal{K}_{2}}V\left( \rho \otimes \xi \right) V^{\ast }.
\tag{2}
\end{equation}

The generalized beam splitting  $  \Pi ^{\ast } $   with the vacuum noise state  $  %
\xi _{0}=\left\vert 0\right\rangle \left\langle 0\right\vert  $   is called the
beam splitter  $  \Pi _{0}^{\ast } $   given by 
\[
\Pi _{0}^{\ast }\left( \left\vert \theta \right\rangle \left\langle \theta
\right\vert \otimes \xi _{0}\right) =\left\vert \alpha \theta \right\rangle
\left\langle \alpha \theta \right\vert \otimes \left\vert-\bar{\beta}\theta
\right\rangle \left\langle-\bar{\beta}\theta \right\vert 
\]%
for the coherent input state  $  \rho \otimes \xi _{0}=\left\vert \theta
\right\rangle \left\langle \theta \right\vert \otimes \left\vert
0\right\rangle \left\langle 0\right\vert \in \mathfrak{S}\left( \mathcal{H}%
_{1}{\otimes \mathcal{K}}_{1}\right)  $. The beam splitter  $  \Pi _{0}^{\ast } $  
was described by means of the lifting  $  \mathcal{E}_{0}^{\ast } $   from  $  %
\mathfrak{S}\left( \mathcal{H}\right)  $   to  $  \mathfrak{S}\left( \mathcal{H}{%
\otimes \mathcal{K}}\right)  $   in the sense of Accardi and Ohya \cite{nv:-AL} as
follows 
\[
\mathcal{E}_{0}^{\ast }\left( \left\vert \theta \right\rangle \left\langle
\theta \right\vert \right) =\left\vert \alpha \theta \right\rangle
\left\langle \alpha \theta \right\vert \otimes \left\vert \beta \theta
\right\rangle \left\langle \beta \theta \right\vert. 
\]%
Based on the liftings, the beam splitting was studied by Accardi--Ohya and
Ficht\-ner--Freudenberg--Libsher \cite{nv:-FFL}.   Moreover, the noisy quantum
channel  $  \Lambda _{0}^{\ast } $   with the vacuum noise state  $  \xi
_{0}=\left\vert 0\right\rangle \left\langle 0\right\vert  $   is called the
attenuation channel given by Ohya \cite{nv:-O1} as 
\begin{equation}
\Lambda _{0}^{\ast }(\rho )\equiv \mathop{\rm tr}\nolimits_{\mathcal{K}_{2}}\Pi _{0}^{\ast }(\rho
\otimes \xi _{0})=\mathop{\rm tr}\nolimits_{\mathcal{K}_{2}}V_{0}(\rho \otimes |0\rangle \langle
0|)V_{0}^{\ast },
\tag{3}
\end{equation}%
which plays an important role for investigating the quantum communication
processes.

\smallskip

\textbf{2. Complexities}

Two kind of complexities  $  C^{\mathcal{S}}( \rho)  $,    $  T^{\mathcal{%
S}}( \rho ; \Lambda ^{\ast } )  $   are used in ID.  $  C^{\mathcal{S}%
}( \rho)  $   is a comple\-xity of a state  $  \rho  $   measured from a
subset  $  \mathcal{S} $   and  $  T^{\mathcal{S}}\left( \rho {;\Lambda ^{\ast }}%
\right)  $   is a transmitted complexity according to the state change from  $  %
\rho  $   to  $  \Lambda ^{\ast }\rho  $. These complexities should fulfill the
following conditions:   Let  $  \mathcal{S}, $    $  \overline{\mathcal{S}} $,    $  \mathcal{S%
}_{t} $   be subsets of  $  \mathfrak{S}\left( \mathcal{H}_{1}\right),  $    $  %
\mathfrak{S}\left( \mathcal{H}_{2}\right),  $    $  \mathfrak{S}\left( \mathcal{H}%
_{1}\otimes \mathcal{H}_{2}\right) $,   respectively.

\begin{itemize}
\item[{(1)}] For any  $  \rho \in \mathcal{S} $,    $  C^{\mathcal{S}}(\rho ) $   and  $  T^{%
\mathcal{S}}( \rho ; \Lambda ^{\ast } )  $   are nonnegative ( $  C^{%
\mathcal{S}}(\rho )\geq 0 $,   \linebreak  $  T^{\mathcal{S}}\left( \rho {;\Lambda ^{\ast }%
}\right) \geq 0 $  ).

\item[{(2)}] For a bijection  $  j $   from  $  \mathop{\rm ex}\mathfrak{S}\left( \mathcal{H}_{1}\right)  $  
to  $  \mathop{\rm ex}\mathfrak{S}\left( \mathcal{H}_{1}\right)  $,   
\[
C^{\mathcal{S}}(\rho )=C^{\mathcal{S}}(j( \rho) ) 
\]%
is hold, where  $  \mathop{\rm ex}\mathfrak{S}\left( \mathcal{H}_{1}\right)  $   is the set of
extremal point of  $  \mathfrak{S}\left( \mathcal{H}_{1}\right)  $.

\item[{(3)}] For  $  \rho \otimes \sigma \in \mathfrak{S}\left( \mathcal{H}_{1}\otimes 
\mathcal{H}_{2}\right)  $,    $  \rho \in \mathfrak{S}\left( \mathcal{H}%
_{1}\right)  $,    $  \sigma \in \mathfrak{S}\left( \mathcal{H}_{2}\right)  $,   
\[
C^{\mathcal{S}_{t}}(\rho \otimes \sigma )=C^{\mathcal{S}}(\rho )+C^{%
\overline{\mathcal{S}}}(\sigma ).
\]%
It means that the complexity of the state  $  \rho \otimes \sigma  $   of totally
independent systems are given by the sum of the complexities of the states  $  %
\rho  $   and  $  \sigma  $.

\item[{(4)}]  $  C^{\mathcal{S}}( \rho)  $   and  $  T^{\mathcal{S}}\left( \rho {%
;\Lambda ^{\ast }}\right)  $   satisfy the following inequality  $  0\leq T^{%
\mathcal{S}}\left( {\rho ;\Lambda ^{\ast }}\right) \leq C^{\mathcal{S}}(\rho
) $.

\item[{(5)}] If the channel  $  {\Lambda ^{\ast }} $   is given by the identity map  $  {id} $,   then  $  T^{\mathcal{S}}\left( {\rho ;id}\right) =C^{\mathcal{S}}({\rho }) $  
is hold.


\end{itemize}

One of the example of the above complexities are the Shannon entropy  $  %
S\left( p\right)  $   for  $  C^{\mathcal{S}}(p) $   and classical mutual entropy  $  %
I\left( p;{\Lambda ^{\ast }}\right)  $   for  $  T^{\mathcal{S}}\left( p{;\Lambda
^{\ast }}\right)  $.   Let us consider these complexities for quantum systems.

\smallskip

\textbf{2.1. Example of Complexity  $  \boldsymbol{C^{\mathcal{S}}\left( \protect\rho %
\right) } $  }

\textit{2.1.1. von Neumann Entropy and  \bf  $  \boldsymbol{\mathcal{S}} $-mixing entropy}

One of the example of the complexity  $  C^{\mathcal{S}}({\rho }) $   of ID in
quantum system is the von Neumann entropy  $  {S( \rho) } $   \cite{nv:-Neu}  described by 
\[
C^{\mathcal{S}}( \rho) \Leftrightarrow {S( \rho)
=-\mathop{\rm tr}\rho \log \rho } 
\]%
for any density operators  $  \rho \in \mathfrak{S}\left( \mathcal{H}%
_{1}\right)  $,   which satisfies the above conditions (1), (2), (3).   

%\hear

Let  $  (\mathcal{A},\mathfrak{S}(\mathcal{A}),\alpha (G)) $   be a C*-dynamical
system and  $  \mathcal{S} $   be a weak* compact and convex subset of  $  \mathfrak{S%
}(\mathcal{A}) $. For example,  $  \mathcal{S} $   is given by  $  \mathfrak{S}(%
\mathcal{A}) $   (the set of all states on  $  \mathcal{A} $),  $  I(\alpha ) $   (the
set of all invariant states for  $  \alpha  $),  $  K(\alpha ) $   (the set of all
KMS states), and so on. Every state  $  \varphi \in \mathcal{S} $   has a maximal
measure  $  \mu  $   pseudosupported on  $  ex\mathcal{S} $   such that 
\begin{equation}
\varphi =\int_{\mathcal{S}}\omega d\mu, 
\tag{4}
\end{equation}%
where  $  \mathop{\rm ex}\mathcal{S} $   is \noindent the set of all extreme points of  $  %
\mathcal{S} $.   The measure  $  \mu  $   giving the above decomposition is not
unique unless  $  \mathcal{S} $   is a Choquet simplex.   We denote the set of all
such measures by  $  M_{\varphi }(\mathcal{S}) $,   and define 
\begin{gather}
D_{\varphi }(\mathcal{S}) =\Bigl\{ M_{\varphi }(\mathcal{S});\quad \exists
\mu _{k}\subset \mathbb{R}^{+}\hbox{ and }\left\{ \varphi _{k}\right\}
\subset \mathop{\rm ex}S \\
s.t.\quad \sum_{k}\mu _{k}=1,\quad \mu =\sum_{k}\mu _{k}\delta
\left( \varphi _{k}\right) \Bigr\}, 
\tag{5}
\end{gather}%
\noindent where  $  \delta (\varphi ) $   is the Dirac measure concentrated on an
initial state  $  \varphi  $. For a measure  $  \mu \in D_{\varphi }(\mathcal{S}) $,  
we put 
\begin{equation}
H(\mu )=-\sum_{k}\mu _{k}\log \mu _{k}.
\tag{6}
\end{equation}%
\noindent The C*-entropy of a state   $  \varphi \in \mathcal{S} $  
with respect to  $  \mathcal{S} $   (\textbf{$  \boldsymbol{\mathcal{S}} $-mixing entropy}) is
defined by 
\begin{equation}
C^{\mathcal{S}}\left( \varphi \right) \Leftrightarrow S^{\mathcal{S}%
}(\varphi )=\left\{ 
\begin{array}{c}
\inf \left\{ H\left( \mu \right) ;\quad \mu \in D_{\varphi }(\mathcal{S}%
)\right\}    \\
+\infty \quad \qquad \quad \hbox{if }D_{\varphi }(\mathcal{S})=\varnothing.%
\end{array}%
\right.  
\tag{7}
\end{equation}%
It describes the amount of information of the state  $  \varphi  $   measured from
the subsystem  $  \mathcal{S} $. We denote  $  S^{\mathfrak{S}(\mathcal{A}%
)}(\varphi ) $   by  $  S(\varphi ) $   if  $  \mathcal{S}=\mathfrak{S}(\mathcal{A}) $. 
It is an extension of von Neumann's entropy.

This entropy (mixing  $  \mathcal{S} $-entropy) of a general state  $  \varphi  $  
satisfies the following properties \cite{nv:-O2}.  

%\begin{theorem}
{\small \sc Theorem 2.1.}{\it
 When  $  \mathcal{A}=B(\mathcal{H}) $   and  $  \alpha _{t}=Ad(U_{t}) $   (i.e.,  $  %
\alpha _{t}(A)=U_{t}^{\ast }AU_{t} $   for any  $  A\in \mathcal{A} $) with a
unitary operator  $  U_{t} $,   for any state  $  \varphi  $   given by  $  \varphi ({}\cdot{}
)=\mathop{\rm tr}\rho {}\cdot {} $   with a density operator  $  \rho  $,   the following
facts hold: 
%\end{theorem}

\begin{enumerate}
\item[\rm (1)]  $  S(\varphi )=-\mathop{\rm tr}\rho \log \rho  $.

\item[\rm (2)] If  $  \varphi  $   is an  $  \alpha  $-invariant faithful state and every
eigenvalue of  $  \rho  $   is non-degenerate, then  $  S^{I(\alpha )}(\varphi
)=S(\varphi ), $   where  $  I( \alpha)  $   is the set of all  $  \alpha  $  %
-invariant faithful states.

\item[\rm (3)] If  $  \varphi \in K(\alpha ) $,   then  $  S^{K( \alpha)
}(\varphi )=0 $,   where  $  K $    $  ( \alpha)  $  is the set of all KMS
states.
\end{enumerate}
}


%\begin{theorem}
{\small \sc Theorem 2.2.}{\it
For any  $  \varphi \in K(\alpha ) $,   we have

\begin{enumerate}
\item[\rm (1)]  $  S^{K( \alpha) }(\varphi )\leq S^{I(\alpha )}(\varphi
) $.

\item[\rm (2)]  $  S^{K( \alpha) }(\varphi )\leq S(\varphi ) $.
\end{enumerate}
}
%\end{theorem}

\smallskip
\textbf{2.2. Example of Transmitted Complexity  $  \boldsymbol{T^{\mathcal{S}}\left( \protect%
\rho ;\Lambda ^{\ast }\right) } $  }

The classical mutual entropy  $  I\left( p;{\Lambda ^{\ast }}\right)  $   defined
by using the joint probability distribution between the input state and the
output state is an example of the transmitted complexity  $  T^{\mathcal{S}%
}\left( p{;\Lambda ^{\ast }}\right)  $   of ID. In general, there does not exit
the joint states in the quantum system \cite{nv:-Urb}.  We need to introduce
the compound state in quantum system instead of the joint probability
distribution in classical system.

\smallskip
\textbf{2.3. Compound state}

The quantum mutual entropy  $  I\left( {\rho, \Lambda ^{\ast }}\right)  $   should
satisfy the following three con\-ditions:

\begin{enumerate}
\item If the channel is given by the identity channel  $  id $,   then  $  I\left(
\rho ;id\right) =S( \rho)  $   (von Neumann entropy) is hold.

\item If the system is classical, then the quantum mutual equals to the
classical mutual entropy.

\item The quantum mutual entropy should satisfy the Shannon's type
inequalities:%
\[
0\leq I\left( {\rho, \Lambda ^{\ast }}\right) \leq {S( \rho) }. 
\]
\end{enumerate}

Ohya introduced two compound states  $  \sigma _{0} $   and  $  \sigma _{E} $.    $  %
\sigma _{0} $   is the trivial compound state given by 
\[
\sigma _{0}=\rho \otimes {\Lambda ^{\ast }\rho.} 
\]%
 $  \sigma _{E} $   is the compound state representing a certain correlation
between the input state and the output state given by%
\[
\sigma _{E}=\sum_{n}\lambda _{n}E_{n}\otimes \Lambda ^{\ast }E_{n} 
\]
associated with the Schatten--von Neumann (one dimensional spectral)
decompo\-sition \cite{nv:-Sch}  $  \rho =\sum\nolimits_{n}\lambda _{n}E_{n} $   of the
input state  $  \rho  $.

\smallskip
\textbf{2.4. Ohya Mutual Entropy for density operator}

An example of the transmitted complexity  $  T^{\mathcal{S}}\left( \rho ; \Lambda ^{\ast }\right)  $   of ID in quantum system is the Ohya mutual
entropy with respect to the initial state  $  \rho  $   and the quantum channel  $  %
\Lambda ^{\ast } $   defined by%
\begin{equation}
T^{\mathcal{S}}\left( {\rho ;\Lambda ^{\ast }}\right) \Leftrightarrow
I\left( \rho ;\Lambda ^{\ast }\right) \equiv \sup \left\{ \sum_{n}S(\sigma
_{E},\sigma _{0}),\rho =\sum_{n}\lambda _{n}E_{n}\right\},  \nonumber
\end{equation}%
where  $  S\left( {}\cdot {}, {}\cdot {}\right)  $   is the Umegaki's relative entropy 
\cite{nv:-Ume}  denoted by 
\begin{equation}
\tag{8}
S\left( {\rho, \sigma }\right) \equiv \left\{ 
\begin{array}{cc}
\mathop{\rm tr}\rho \left( \log \rho-\log \sigma \right) & \left( \hbox{when }\overline{%
\mathop{\rm ran}\rho }\subset \overline{\mathop{\rm ran}\sigma }\right)    \\
\infty & \left( \hbox{otherwise}\right)%
\end{array}%
\right.
\end{equation}%
which was extended to more general quantum systems by Araki and Uhlmann \cite{nv:-O3,nv:-OP,nv:-OV, nv:-Ara,nv:-Uhl}. The Ohya mutual entropy holds the above conditions
(1), (4), (5) such as%
\[
0\leq I\left( {\rho, \Lambda ^{\ast }}\right) \leq {S( \rho) }, 
\]%
\[
I\left( {\rho, id}\right) ={S( \rho).} 
\]%
The capacity means the ability of the information transmission of the
channel, which is used as a measure for construction of channels. The
quantum capacity is formulated by taking the supremum of the Ohya mutual
entropy with respect to a certain subset of the initial state space. The
quantum capacity of quantum channel was studied in [\citen{nv:-OPW1}--\citen{nv:-OW2}].

%\begin{theorem}
{\small\sc Theorem 2.3.} {\it
Let   $\Phi _{E} $  be  a  compound  
state  w.r.t.  the  initial  state $\rho$,   the  quan\-tum  CP
  channel $  \Lambda ^{\ast }$  {and  a%
  Schatten  decomposition  of } $  {\rho
=\sum_{k}\lambda _{k}E_{k}} $  {  defined  by}%
\[
{\Phi _{E}=\sum_{n}\left( I\otimes V_{n}\right) \left[ \sum_{k}\sqrt{\lambda
_{k}}{\left\vert x_{k}\right\rangle }\otimes {\left\vert x_{k}\right\rangle }%
\right] \left[ \sum_{k^{\prime }}\sqrt{\lambda _{k^{\prime }}}{\left\langle
x_{k^{\prime }}\right\vert }\otimes {\left\langle x_{k^{\prime }}\right\vert 
}\right] \left( I\otimes V_{n}^{\ast }\right) }
\]%
under  the  condition   %
\[
{\sum_{n}\left( I\otimes V_{n}^{\ast }\right) \left( I\otimes V_{n}\right)
=I\otimes I\mathrm{\;}}
\]%
and $   \Lambda ^{\ast }$  is  given  by $   %
\Lambda ^{\ast }( \rho) =\sum_{n}V_{n}\rho V_{n}^{\ast } $. By
defining the compound state  $  {\Phi _{E}} $,    one can obtain the following
theorem.
%\end{theorem}

}

\smallskip

%\begin{theorem}
{\small\sc Theorem 2.4.} {\it
    {For the compound state  } $  {\Phi _{E}} $  given above, one can  obtain  two  marginal  states  as  follows 
\begin{eqnarray*}
{\tr_{{\mathcal{H}}_{2}}\Phi _{E}}{=S( \rho), } && \\
{\tr_{{\mathcal{H}}_{1}}\Phi _{E}}{=S\left( \Lambda ^{\ast }\rho \right).} &&
\end{eqnarray*}%
The upper bound of the relative entropy  $  S{\left( \Phi _{E},\rho \otimes
\Lambda ^{\ast }\rho \right) } $   is obtained as follows:%
\[
S{\left( \Phi _{E},\rho \otimes \Lambda ^{\ast }\rho \right) \leq 2S\left(
\rho \right).}
\]
}

\smallskip

{Let} $ \Psi _{E}$  {be  a  compound%
  state  defined} {by}%
\[
{\Psi _{E,\mu }=\mu \;\sigma _{E}+\left( 1-\mu \right) \Phi _{E}{\quad }%
\left( \mu \in \left[ 0,1\right] \right).} 
\]%
We have the following theorem.


%\begin{theorem}
{\small\sc Theorem 2.5.} {\it
    {For the compound state  } $  {\Psi _{E,\mu }} $  w.r.t.  $  \mu \in %
\left[ 0, 1\right]  $   given above,  one  can  obtain  two  marginal  states as  follows%
\begin{eqnarray*}
{tr_{{\mathcal{H}}_{2}}\Psi _{E,\mu }}{=S( \rho), } && \\
{tr_{{\mathcal{H}}_{1}}\Psi _{E,\mu }}{=S\left( \Lambda ^{\ast }\rho \right)
.} &&
\end{eqnarray*}%
The upper bound of the relative entropy  $  S{\left( \Psi _{E,\mu },\rho
\otimes \Lambda ^{\ast }\rho \right) } $   is obtained as follows:%
\[
S{\left( \Psi _{E,\mu },\rho \otimes \Lambda ^{\ast }\rho \right) \leq
\left( 2-\mu \right) S( \rho) \quad }\left( \mu \in \left[ 0,1%
\right] \right) {.}
\]
%\end{theorem}
}

\smallskip

\textbf{2.5. Ohya Mutual Entropy for general C*-system}

Let  $  \left( {{\mathcal{A}},\mathfrak{S}(\mathcal{A}),\alpha \left( G\right) }%
\right)  $   be a unital  $  C^{\ast } $-system and  $  {\mathcal{S}} $   be a weak*
compact convex subset of  $  \mathfrak{S}(\mathcal{A}) $. For an initial state  $  %
\varphi \in \mathcal{S} $   and a channel  $  \Lambda ^{\ast }\,:\,\mathfrak{S}%
\left( {\mathcal{A}}\right) \rightarrow \mathfrak{S}\left( {\mathcal{B}}%
\right)  $,   two compound states \cite{nv:-OP, nv:-O2} are defined by%
\begin{gather}
\Phi _{\mu }^{\mathcal{S}} =\int_{\mathcal{S}}\omega \otimes \Lambda
^{\ast }\omega \;d\mu, \tag{9} \\
\Phi _{0} =\varphi \otimes \Lambda ^{\ast }\varphi.
\tag{10}
\end{gather}%
\noindent The compound state  $  \Phi _{\mu }^{\mathcal{S}} $   expresses the
correlation between the input state  $  \varphi  $   and the output state  $  \Lambda
^{\ast }\varphi  $. The mutual entropy with respect to  $  \mathcal{S} $   and  $  \mu 
 $   is given by 
\begin{equation}
\tag{11}
I_{\mu }^{\mathcal{S}}\left( {\varphi ;\Lambda ^{\ast }}\right) =S\left( {%
\Phi _{\mu }^{\mathcal{S}},\Phi _{0}}\right)
\end{equation}%
\noindent and the mutual entropy with respect to  $  \mathcal{S} $   is defined by
Ohya \cite{nv:-OP, nv:-O2}  as 
\begin{equation}
T^{\mathcal{S}}\left( {\varphi ;\Lambda ^{\ast }}\right) \Leftrightarrow I^{%
\mathcal{S}}\left( {\varphi ;\Lambda ^{\ast }}\right) =\sup \left\{ {I_{\mu
}^{\mathcal{S}}\left( {\varphi ;\Lambda ^{\ast }}\right) \;;\;\mu \in
M_{\varphi }\left( {\mathcal{S}}\right) }\right\}.
\tag{12}
\end{equation}


\smallskip

\textbf{2.6. Other Mutual Entropy Type Measures}

Recently, several mutual entropy type measures were proposed by Shor \cite{nv:-Sho} and Bennet et al \cite{nv:-BNS,nv:-BSST}, which defined by using the entropy
exchange \cite{nv:-Schu} given by%
\begin{equation}
\tag{13}
S_{e}\left( {\rho, \Lambda ^{\ast }}\right) =-\tr W\log W\hbox{,}
\end{equation}%
where  $  W $   is a matrix  $  W=\left( {W_{ij}}\right) _{i,j} $   with the elements 
\begin{equation}
W_{ij}\equiv \tr A_{i}^{\ast }\rho A_{j}
\tag{14}
\end{equation}%
obtained by means of the input state  $  \rho  $   and the CP channel  $  \Lambda
^{\ast } $   described by a Stinespring--Sudarshan--Kraus form 
\begin{equation}
\Lambda ^{\ast }\left({} \cdot {}\right) \equiv \sum\nolimits_{j}A_{j}^{\ast
}\cdot A_{j}.
\tag{15}
\end{equation}%
Based on the entropy exchange, the coherent entropy  $  I_{C}\left( {\rho
;\Lambda ^{\ast }}\right)  $   \cite{nv:-SN} and the Lindblad-Nielson entropy  $  %
I_{L}\left( {\rho ;\Lambda ^{\ast }}\right)  $   \cite{nv:-BSST} were defined by%
\begin{gather}
I_{C}\left( {\rho ;\Lambda ^{\ast }}\right) \equiv S\left( {\Lambda ^{\ast
}\rho }\right)-S_{e}\left( {\rho, \Lambda ^{\ast }}\right), 
\tag{16}
\\
I_{L}\left( {\rho ;\Lambda ^{\ast }}\right) \equiv S( \rho)
+S\left( {\Lambda ^{\ast }\rho }\right)-S_{e}\left( {\rho, \Lambda ^{\ast }}%
\right).
\tag{17}
\end{gather}

\smallskip

\textbf{2.7. Comparison among these quantum mutual entropy type measures}

In this section, we compare with these mutual types measures.

By comparing these mutual entropies for quantum information communication
processes, we have the following theorem \cite{nv:-OW4}:

%\begin{theorem}

\smallskip
{\small\sc Theorem 2.6.} {\it 
Let  $  \left\{ {A_{j}}\right\}  $   be a projection valued measure with dim $  %
A_{j}=1 $. For arbitrary state  $  \rho  $   and the quantum channel  $  \Lambda
^{\ast }\left( {}\cdot {}\right) \equiv \sum\nolimits_{j}A_{j}\cdot A_{j}^{\ast
}  $,   one has

\begin{itemize}

\item[ $  (1) $]    $  0\leq I\left( {\rho ;\Lambda ^{\ast }}\right) \leq \min \left\{ {%
S( \rho), S\left( {\Lambda ^{\ast }\rho }\right) }\right\}  $  %
(Ohya mutual entropy),

\item[ $  (2) $]    $  I_{C}\left( {\rho ;\Lambda ^{\ast }}\right) =0 $   (coherent entropy),

\item[ $  (3) $]    $  I_{L}\left( {\rho ;\Lambda ^{\ast }}\right) =S( \rho)  $  
(Lindblad entropy).

\end{itemize}
%\end{theorem}
}

\smallskip

For the attenuation channel  $  \Lambda _{0}^{\ast } $,   one can obtain the
following theorems \cite{nv:-OW4}:

\smallskip
%\begin{lemma}
{\small \sc Lemma 2.1.}
{\it 
For  the  attenuation  channel $  \Lambda _{0}^{\ast }$  %
and  the  input  state%
\[
{\rho =\lambda {\left\vert 0\right\rangle }{\left\langle 0\right\vert }%
+(1-\lambda ){\left\vert \theta \right\rangle }{\left\langle \theta
\right\vert },} 
\]%
there  exists  a  unitary  operator $  U$  such  that%
\[
{UWU^{\ast }=\lambda {\left\vert 0\right\rangle }{\left\langle 0\right\vert }%
+(1-\lambda ){\left\vert-\overline{\beta }\theta \right\rangle }{%
\left\langle-\overline{\beta }\theta \right\vert }.} 
\]

}
%\end{lemma}

\smallskip

%\begin{theorem}
{\small\sc Theorem 2.7.}
{\it 
For  the  attenuation  channel $  \Lambda _{0}^{\ast }
 $  and  the  input  state %
\[
{\rho =\lambda {\left\vert 0\right\rangle }{\left\langle 0\right\vert }%
+(1-\lambda ){\left\vert \theta \right\rangle }{\left\langle \theta
\right\vert },} 
\]%
the  entropy  exchange  is  obtained  by  %
\[
{S_{e}\left( \rho, \Lambda _{0}^{\ast }\right) =-\tr W\log
W=-\sum_{j=0}^{1}\mu _{j}\log \mu _{j},} 
\]%
where%
\[
{\mu _{j}=\frac{1}{2}\left\{ 1+\left(-1\right) ^{j}\sqrt{1-4\lambda \left(
1-\lambda \right) \left( 1-\exp \left(-\left\vert \beta \right\vert
^{2}\left\vert \theta \right\vert ^{2}\right) \right) }\right\} \quad \left(
j=0, 1\right).} 
\]
}
%\end{theorem}

\smallskip

%\begin{theorem}
{\small\sc Theorem 2.8.}
{\it 
For any state  $  \rho =\sum_{n}{\lambda _{n}}\left\vert n\right\rangle
\left\langle n\right\vert  $   and the attenuation channel  $  \Lambda _{0}^{\ast
}  $   with  $  \left\vert \alpha \right\vert ^{2}=\left\vert \beta \right\vert
^{2}=\frac{1}{2} $,   one has

\begin{itemize}

\item[$  (1) $]    $  0\leq I\left( {\rho ;\Lambda _{0}^{\ast }}\right) \leq \min \left\{ {%
S( \rho), S\left( {\Lambda _{0}^{\ast }\rho }\right) }\right\}  $  %
(Ohya mutual entropy),

\item[$  (2) $]    $  I_{C}\left( {\rho ;\Lambda _{0}^{\ast }}\right) =0 $   (coherent
entropy),

\item[$  (3) $]   $  I_{L}\left( {\rho ;\Lambda _{0}^{\ast }}\right) =S( \rho) 
 $   (Lindblad entropy).

\end{itemize} 
 
%\end{theorem}
}
\smallskip

%\begin{theorem}
{\small\sc Theorem 2.9.}
{\it 
For the attenuation channel  $  \Lambda _{0}^{\ast } $   and the input state  $$  \rho
=\lambda \left\vert 0\right\rangle \left\langle 0\right\vert +(1-\lambda
)\left\vert \theta \right\rangle \left\langle \theta \right\vert  ,$$   we have
\begin{itemize}
\item[$  (1) $]    $  0\leq I\left( {\rho ;\Lambda _{0}^{\ast }}\right) \leq \min \left\{ {%
S( \rho), S\left( {\Lambda _{0}^{\ast }\rho }\right) }\right\}  $  %
(Ohya mutual entropy),

\item[$  (2) $]    $-S( \rho) \leq I_{C}\left( {\rho ;\Lambda _{0}^{\ast }}%
\right) \leq S( \rho)  $   (coherent entropy),

\item[$  (3) $]    $  0\leq I_{L}\left( {\rho ;\Lambda _{0}^{\ast }}\right) \leq 2S\left(
\rho \right)  $   (Lindblad entropy).

\end{itemize}
%\end{theorem}
}

\smallskip

It shows that the coherent entropy holds  $  I_{C}\left( {\rho ;\Lambda
_{0}^{\ast }}\right)  $    $  <0 $   for  $  \left\vert \alpha \right\vert
^{2}<\left\vert \beta \right\vert ^{2} $   and the Lindblad entropy satisfies  $  %
I_{L}\left( {\rho ;\Lambda _{0}^{\ast }}\right) \geq S( \rho)  $  
for  $  \left\vert \alpha \right\vert ^{2}>\left\vert \beta \right\vert ^{2} $.
From the above theorems, we can conclude that the transmitted complexity in
quantum system is the Ohya mutual entropy and it is most fitting measure for
studying the efficiency of information transmission in quantum communication
processes. It means that Ohya mutual entropy can be considered as the
transmitted complexity for quantum communication processes.

\smallskip

%\begin{theorem}
{\small\sc Theorem 2.10.}
{\it 
    For  the  attenuation  channel  $  \Lambda _{0}^{\ast }$  and  the  input
state%
\[
{\rho =\lambda {\left\vert 0\right\rangle }{\left\langle 0\right\vert }%
+(1-\lambda ){\left\vert \theta \right\rangle }{\left\langle \theta
\right\vert },} 
\]
if $  \lambda =\frac{1}{2} $  and $  \beta =\sqrt{\frac{2}{3}} $,  then  there  exists  a
compound  state $  \Phi $  satisfying%
\[
I_{L}{\left( \rho ;\Lambda _{0}^{\ast }\right) =S\left( \Phi, \rho \otimes
\Lambda _{0}^{\ast }\rho \right) \mathrm{.}} 
\]
%\end{theorem}
}

\smallskip

%\begin{theorem}
{\small\sc Theorem 2.11.}
{\it 
For  the  attenuation  channel  $  \Lambda _{0}^{\ast } $  and  the  input
state%
\[
{\rho =\lambda {\left\vert 0\right\rangle }{\left\langle 0\right\vert }%
+(1-\lambda ){\left\vert \theta \right\rangle }{\left\langle \theta
\right\vert },} 
\]
if $   \lambda =\frac{1}{2} $  and $   \alpha =1 $,  then  there  exists  a  compound  state $  \Phi $  satisfying%
\[
{S\left( \Phi, \rho \otimes \Lambda _{0}^{\ast }\rho \right) =S\left( \rho
\right).} 
\]
}



\thanks{
\textbf{Acknowledgment.}
The author would like to thank Prof. I.~Volovich for his many helpful
sugestions during the preparation of the paper.
}


\def\refname{REFERENCES}
\begin{thebibliography}{99}



\Bibitem{nv:-O3} 
\by M.~Ohya
\paper Information dynamics and its application to optical communication processes
\inbook Quantum Aspects of Optical Communications
\serial Lecture Notes in Physics 
\procinfo Paris, 1990
\vol 378
\yr 1991
\publaddr Berlin
\publ Springer
\pages 81--92 
%DOI 10.1007/3-540-53862-3_169



\Bibitem{nv:-O1} 
\by M.~Ohya
\paper On compound state and mutual information in quantum information theory
\jour Information Theory, IEEE Trans. 
\vol 29
\issue 5
\pages 770--774
\yr 1983
%Digital Object Identifier :  10.1109/TIT.1983.1056719



\Bibitem{nv:-OP} 
\by M.~Ohya, D.~Petz
\book  Quantum entropy and its use
\serial Texts and Monographs in Physics
\publaddr Springer Verlag
\publ Berlin
\yr 1993
\totalpages viii+335


\Bibitem{nv:-OV} 
\by M.~Ohya, I.~Volovich
\serial Theoretical and Mathematical Physics
\yr 2011
\book Mathematical foundations of quantum information and computation and its applications to nano- and bio-systems
\totalpages xx+759
\publ Springer 
\publaddr Dordrecht
%DOI 10.1007/978-94-007-0171-7


\Bibitem{nv:-Ing} 
\by R.~S.~Ingarden, A.~Kossakowski, M.~Ohya
\book Information dynamics and open systems. Classical and quantum approach
\serial Fundamental Theories of Physics
\publ Kluwer Academic Publ.
\vol 86
\publaddr Dordrecht
\yr 1997
\totalpages x+307



\Bibitem{nv:-OW3} 
\by M.~Ohya, N.~Watanabe
\paper  Construction and analysis of a mathematical model in quantum communication processes
\jour Electronics and Communications in Japan, Part~1
\vol 68
\issue 2
\pages 29--34 
\yr 1985
%DOI: 10.1002/ecja.4410680204

\Bibitem{nv:-AL} 
\by L.~Accardi, M.~Ohya
\paper Compound channels, transition expectations, and liftings
\jour Appl. Math. Optim.
\vol 39
\issue 1
\pages 33--59 
\yr 1999
%DOI     10.1007/s002459900097


\Bibitem{nv:-FFL} 
\by K.-H.~Fichtner, W.~Freudenberg, V.~Liebscher
\preprint Beam splittings and time evolutions of Boson systems
\preprintinfo  Forschungsergebnise der Fakult\"at f\"ur Mathematik und Informatik, 
\href{http://www.math-inf.uni-greifswald.de/biomathematik/liebscher/publications/ffl.ps}{Math/Inf/96/39}
\totalpages 105 
\yr 1996


\Bibitem{nv:-Neu} 
\by J.~von Neumann
\book Die Mathematischen Grundlagen der  Quantenmechanik (German) 
\publ Springer Verlag
\publaddr Berlin
\yr 1932
\totalpages v+262


\Bibitem{nv:-O2} 
\by M.~Ohya
\paper Some aspects of quantum information theory and their applications to irreversible processes
\jour Rep. Math. Phys.
\vol 27
\issue 1
\pages 19--47
\yr 1989


\Bibitem{nv:-Urb} 
\by K.~Urbanik
\paper Joint probability distribution of observables in quantum mechanics
\jour Studia Math.
\vol 21
\issue 1
\pages 117--133
\yr 1961


\Bibitem{nv:-Sch} 
\by R.~Schatten
\book  Norm ideals of completely continuous operators
\serial Ergebnisse der Mathematik und ihrer Grenzgebiete
\vol 27
\publ Springer Verlag
\publaddr Berlin, New York
\yr 1970
\totalpages vii+81


\Bibitem{nv:-Ume} 
\by H.~Umegaki
\paper Conditional expectation in an operator algebra. IV. Entropy and information
\jour K\={o}dai Math. Sem. Rep.
\vol 14
\pages 59--85 
\yr 1962


\Bibitem{nv:-Ara} 
\by H.~Araki
\paper Relative entropy for states of von Neumann algebras
\jour Publ. Res. Inst. Math. Sci. 
\vol 11
\issue 3
\pages 809--833
\yr 1976


\Bibitem{nv:-SN} 
\by B.~W.~Schumacher, M.~A.~Nielsen
\paper Quantum data processing and error correction
\jour Phys. Rev.~A
\vol 54
\issue 4
\pages 2629–2635
\yr 1996
\arxiv \href{http://arxiv.org/abs/quant-ph/9604022}{quant-ph/9604022}

\Bibitem{nv:-Uhl} 
\by A.~Uhlmann
\paper Relative entropy and the Wigner--Yanase--Dyson--Lieb concavity in interpolation theory
\jour Commun. Math. Phys.
\vol 54
\issue 1
\pages 21--32
\yr 1977


\Bibitem{nv:-OPW1} 
\by M.~Ohya, D.~Petz, N.~Watanabe
\paper On capacity of quantum channels
\jour Probab. Math. Statist. 
\vol 17
\issue 1
\issueinfo Acta Univ. Wratislav. No.~1928
\pages 179--196 
\yr 1997




\Bibitem{nv:-OPW2} 
\by M.~Ohya, D.~Petz, N.~Watanabe
\paper Numerical computation of quantum capacity
\paperinfo Proc. of the International Quantum Structures Association (Berlin, 1996)
\jour Internat. J. Theoret. Phys.
\vol 37
\issue 1
\pages 507--510 
\yr 1998

 



\Bibitem{nv:-OW1} 
\by M.~Ohya, N.~Watanabe
\paper Quantum capacity of noisy quantum channel
\jour Quantum Communication and Measurement
\vol 3
\pages 213--220 
\yr 1997

\Bibitem{nv:-OW2} 
\by M.~Ohya, N.~Watanabe
\book Foundation of Quantum Communication Theory (in Japanese)
\publ Makino Pub. Co.
\yr 1998

\Bibitem{nv:-Sho} 
\by P.~Shor
\preprint \href{http://www.msri.org/publications/ln/msri/2002/quantumcrypto/shor/1/meta/aux/shor.pdf}{The quantum channel capacity and coherent information}
\preprintinfo Lecture Notes, MSRI Workshop on Quantum Computation
\yr 2002


\Bibitem{nv:-BNS} 
\by H.~Barnum, M.~A.~Nielsen, B.~W.~Schumacher
\paper Information transmission through a noisy quantum channel
\jour Phys. Rev.~A
\vol 57
\issue 6
\pages 4153--4175 
\yr 1998
\arxiv 	\href{http://arxiv.org/abs/quant-ph/9702049}{quant-ph/9702049}

\Bibitem{nv:-BSST} 
\by C.~H.~Bennett, P.~W.~Shor, J.~A.~Smolin, A.~V.~Thapliyalz
\paper Entanglement-assisted capacity of a quantum channel and the reverse Shannon theorem
\arxiv \href{http://arxiv.org/abs/quant-ph/0106052}{quant-ph/0106052}
\yr 2001
\jour Inform. Theory, IEEE Trans.
\vol 48 
\yr 2002
\issue 10
\pages 2637–2655

\Bibitem{nv:-Schu} 
\by B.~W.~Schumacher
\paper Sending entanglement through noisy quantum channels
\jour Phys. Rev.~A
\vol 54
\issue 4
\yr 1996
\pages 2614–2628 



\Bibitem{nv:-OW4} 
\by M.~Ohya, N.~Watanabe
\preprint Comparison of mutual entropy-type measures
\preprintinfo TUS preprint







\end{thebibliography}


\lfoot[]{}
\rfoot[]{}


%\thispagestyle{empty}

\vskip1mm \par{\footnotesize
\hskip6.5cm \theengdate \par
\hskip6.5cm\theengrevdate\par}



\newpage

%
%\end{document}
